% DOCUMENT CLASS %
\documentclass[a4paper, twoside, 12pt]{book}   % for A4 Paper, 12pt fontsize.
\setcounter{tocdepth}{4}
\setcounter{secnumdepth}{3}

% PACKAGES%
\usepackage{amsfonts}   % For a basic mathfont (many will be replaced by stix)
\usepackage{amsmath}    % For basic math symbols
%\usepackage{bbm}        % For \mathbbm (better version of \mathbb)
\usepackage{mathrsfs}   % For \mathscr
\usepackage{hyperref}   % For footnotes
\usepackage{csquotes}   % For \textquote{}
\usepackage{IEEEtrantools}  % For better alignments
\usepackage{tikz}   % For a macro
\usepackage{listings} % For code
\usepackage{xcolor}   % For ... code
\usepackage[stretch=40]{microtype}
\usepackage{enumitem}
\usepackage{geometry}
\usepackage[color=gray]{todonotes}

% Language
\usepackage[ngerman]{babel}

% If needed
\usepackage[nohyperlinks,smaller]{acronym}

% For tests
\usepackage{lipsum}

%%% ENVIRONMENTS %%%
\usepackage{amsthm}
\usepackage[nameinlink]{cleveref}

% Definition environment
\theoremstyle{definition}
\newtheorem{definition}{Definition}[section]
% Configure reference name
\crefname{definition}{Definition}{Definitions}

\theoremstyle{plain}
\newtheorem{problem}{Problem}%[section]
\crefname{problem}{Problem}{Problems}

\theoremstyle{definition}
\newtheorem{algorithm}{Algorithm}%[section]
\crefname{algorithm}{Algorithm}{Algorithms}

\theoremstyle{theorem}
\newtheorem{theorem}{Theorem}%[section]
\crefname{theorem}{Theorem}{Theorems}
%%%

%%% FONT CONFIGURATION %%%
\usepackage{fontspec}

% Updated version of Crimson Text
\setmainfont{Cochineal}
\setmonofont{Cascadia Code}[Scale=0.8]
\usepackage[cochineal]{newtxmath} % To use Cochineal as a Math Font
%%%

% MAKE DESCRIPTION ENVIRONMENTS HAVE ITALICIZED ITEMS %
\renewcommand{\descriptionlabel}[1]{\hspace{\labelsep}\textit{#1}}

% SETS %
\newcommand*{\R}{\mathbb R}    % Set of real numbers
\newcommand*{\N}{\mathbb N}    % Set of natural numbers
\newcommand*{\Z}{\mathbb Z}    % Set of integers
\newcommand*{\Q}{\mathbb Q}    % Set of rational numbers

% MACROS %
\newcommand*{\QDp}{{:}\text{ }} % Macro for ": " so that when writing something like "\forall x:" there is not a seperation between the '\forall' and the ':'
\newcommand*{\gap}{\text{ }}    % Macro for a lazy aligment 2
\newcommand*{\qedbox}{\null\nobreak\hfill\ensuremath{\square}} % Macro for a QED-Box
% for a funny ¯\_(ツ)_/¯-Emoji
\newcommand{\shrug}[1][]{%
    \begin{tikzpicture}[baseline,x=0.8\ht\strutbox,y=0.8\ht\strutbox,line width=0.125ex,#1]
        \def\arm{(-2.5,0.95) to (-2,0.95) (-1.9,1) to (-1.5,0) (-1.35,0) to (-0.8,0)};
        \draw \arm;
        \draw[xscale=-1] \arm;
        \def\headpart{(0.6,0) arc[start angle=-40, end angle=40,x radius=0.6,y radius=0.8]};
        \draw \headpart;
        \draw[xscale=-1] \headpart;
        \def\eye{(-0.075,0.15) .. controls (0.02,0) .. (0.075,-0.15)};
        \draw[shift={(-0.3,0.8)}] \eye;
        \draw[shift={(0,0.85)}] \eye;
        % draw mouth
        \draw (-0.1,0.2) to [out=15,in=-100] (0.4,0.95);
    \end{tikzpicture}
}
\newcommand{\splitpage}[2]{
    \begin{minipage}{0.5\textwidth}
        #1
    \end{minipage}
    \hspace{0.5cm}
    \begin{minipage}{0.5\textwidth}
        #2
    \end{minipage}
}

% BASIC CONFIGURATION %
\pagestyle{plain}
\allowdisplaybreaks

% CODE %
\definecolor{codepurple}{HTML}{7f0055}
\definecolor{comment}{RGB}{0,128,0}
\definecolor{number}{RGB}{9,136,90}
\definecolor{string}{RGB}{163,21,21}

\lstdefinestyle{default}{
    language=c,
    basicstyle=\ttfamily\footnotesize,  % the size of the fonts that are used for the code
    breakatwhitespace=false,            % sets if automatic breaks should only happen at whitespace
    breaklines=true,                    % sets automatic line breaking
    frame=leftline,
    keepspaces=true,                    % keeps spaces in text, useful for keeping indentation of code (possibly needs columns=flexible)
    numbers=left,                       % where to put the line-numbers; possible values are (none, left, right)
    firstnumber=1,
    numbersep=5pt,                      % how far the line-numbers are from the code
    showspaces=false,                   % show spaces everywhere adding particular underscores; it overrides 'showstringspaces'
    showstringspaces=false,             % underline spaces within strings only
    showtabs=false,                     % show tabs within strings adding particular underscores
    tabsize=2,                          % sets default tabsize to 4 spaces
    morekeywords={__m256, __m256d, __m256i, omp_lock_t, MPI_Comm, MPI_Datatype, MPI_Status},
    keywordstyle=\bfseries\color{codepurple},
    stringstyle=\color{string},
    commentstyle=\itshape\color{comment},
    morecomment=[f]{\#pragma},
    texcl=true,
    mathescape=true,
    xleftmargin=.1\textwidth, xrightmargin=.1\textwidth
}

\lstset{style=default}
